% Refine from the long post on \texttt{10B} on Forum + other stuff 

\section{\texttt{Challenge 10B}}

If this is your first year, \texttt{Challenge 10B} is the final and most difficult challenge of the season. If you've done this before, you already know the deal. 

One of the most difficult things I find with \texttt{10B} every year, apart from the cipher itself, is knowing when to take a break. 

In the midst of chaotic, spiralling ideas on how to crack the final cipher, it’s so easy to fall into that trap of `I’m so close… just a bit longer…' and then 5 hours later your eyes hurt, brain’s fried, and you’re resisting the urge to throw your laptop out the window. 

It feels completely counterintuitive, but stepping away really does help. And even knowing that, I still find it hard to move away when I'm in the zone. 

Over time, I've found a few things that helped me reset and come back with a clearer head. None of these solve the cipher, but they make it much easier to think when you return: 

\textbf{Write Down Everything You Know}\\
Every hour or just pick a timeframe, jot down everything you’ve figured out about the cipher, even the obvious stuff, especially the obvious stuff. Key lengths, repeated n-grams, letter positions, ciphertext arrangement, possible cribs, structure. 

It’s easy to feel like you’re not making much progress, but when you see it on paper, it’s surprising how much you’ve actually uncovered, and it helps track progress. 

\textbf{Have a snack}\\
The brain burns a ridiculous amount of energy when you’re problem-solving. As I’ve been told many times in training: food is fuel.

\textbf{Check in with your team if you have one}\\
It's easy to disappear into your own rabbit hole of ideas so it’s always useful take a moment to sync and share weird hunches, and remind each other to take breaks as well.

\textbf{Setting Expectations}\\
Setting expectations at home if need be. My family is always made well aware that I will be busy during cipher challenge season, especially towards the end so they know what to or not to expect from me. I find that this definitely does reduce some background stress 

\textbf{Stepping Away}\\
Slightly unconventional, but I have a hoodie I only wear when working on the cipher challenge\footnote{Sadly it doesn’t have any actual powers, but having a dedicated codebreaking hoodie does make me feel extra geeky}. 

When I take a break, I physically change out of it. It sounds silly but that very act of changing the hoodie helps me switch mental gears and actually rest. As much as we all love the challenge, I’m sure we all have other things (including food and sleep) that need getting done and I just found that having a physical reminder helped me compartmentalise

\textbf{Sleep}\\
I don't think I need to explain this further. Sleep is good for you, so please remember to sleep at a somewhat sensible time.


\section{Sticking With It}

Having done the cipher challenge for five years, my only real regret was not sticking with \texttt{10B} till the end in my second year. 

I gave up too early and convinced I was going nowhere after working on it for just over a day, only to see others break through with just a bit more persistence. 

What I’ve learned since then is that hints do come, but you’ve got to have patience. They’re mostly be subtle or buried under obscure forum comments, or phrased in a way that only makes sense after you’ve cracked it. 

\section{Celebrating the Little Wins}

It's very easy to let the whole thing feel overwhelming. So try and celebrate the little victories. Remember that marginal gains eventually add up. 

Yes, winning is great, but it's not all about the winning. It sounds cliché to say \textit{value the process over the outcome} but it's actually true. With any performance comes pressure. You start wanting to do well, to keep up and get things right. And slowly, without realising, that pressure can suck the joy out of the experience. 

Don't let the desire to do well take away the reason you started in the first place; don't let the need to perform rip off the curiosity, enjoyment and satisfaction of figuring something out for yourself. 

It's quite funny but I think that if you’re annoyed at a cipher, that’s kind of amazing in itself. 

It means you’ve found something that you care about. Something that's gripped your attention so much that you're willing to struggle with and keep thinking even when it would be easier to stop, and that's rare. 

Most people go through life bored by what they do because let’s be real, if a bunch of weirdly spaced letters and numbers are making you want to throw your keyboard across the room, maybe that’s just be a sign you’ve found something worth fighting with (and for), and that’s a beautiful thing. 

Remember that the feeling of finally solving it is always worth the storm of frustration and patience and mild existential crisis or whatever mix of emotions it takes to get there. 

\section{More than Just Codebreaking}

If anything, this challenge teaches you more than just codebreaking. 

Over the holidays especially, cipher challenge has been escapism in the best sense for me. There's something deeply comforting about knowing that no matter what else is going on, there's always a cipher waiting to be cracked, and a community engaged in doing the same.

And in many ways, the lessons you learn during the competition stay with you as well. Some of the most difficult things I’ve faced over the past few months were made more manageable because of what I learned while struggling with the challenges. Patience, persistence, and the willingness to keep trying different ideas even when the solution isn’t obvious are skills that apply far beyond cryptography. 
